\section{Билет 23. Общее волновое уравнение. Скорость волны в тонком стержне с выводом. Модуль Юнга. Скорость поперечных волн в изотропной твердой среде.}

\begin{definition}
\textbf{Общее волновое уравнение} в трёхмерном пространстве для скалярного возмущения $\xi(\vec{r}, t)$ имеет вид:
\begin{equation}
    \nabla^2 \xi = \frac{1}{v^2} \frac{\partial^2 \xi}{\partial t^2}, \label{eq:general_wave_23}
\end{equation}
где $\nabla^2 = \frac{\partial^2}{\partial x^2} + \frac{\partial^2}{\partial y^2} + \frac{\partial^2}{\partial z^2}$ --- оператор Лапласа, $v$ --- фазовая скорость распространения волны в среде. Это уравнение описывает распространение малых упругих возмущений в однородной изотропной среде без дисперсии и затухания.
\end{definition}

\begin{theorem}[Скорость продольной волны в тонком стержне]
Рассмотрим тонкий упругий стержень с площадью поперечного сечения $S$, плотностью $\rho$ и модулем Юнга $E$. Выделим элементарный объём длиной $dx$. При распространении продольной волны частицы стержня смещаются вдоль оси $x$ на величину $\xi(x,t)$.

Относительная деформация элемента равна $\varepsilon = \partial \xi / \partial x$. Согласно закону Гука, механическое напряжение $\sigma = E \varepsilon = E \, \partial \xi / \partial x$. Разность сил упругости на гранях элемента:
\begin{equation}
    F = S[\sigma(x+dx) - \sigma(x)] = S \frac{\partial \sigma}{\partial x} dx = ES \frac{\partial^2 \xi}{\partial x^2} dx.
\end{equation}
Масса элемента $dm = \rho S dx$. По второму закону Ньютона:
\begin{equation}
    (\rho S dx) \frac{\partial^2 \xi}{\partial t^2} = ES \frac{\partial^2 \xi}{\partial x^2} dx.
\end{equation}
Сокращая на $S dx$, получаем волновое уравнение:
\begin{equation}
    \frac{\partial^2 \xi}{\partial t^2} = \frac{E}{\rho} \frac{\partial^2 \xi}{\partial x^2}.
\end{equation}
Сравнивая с общим видом $\partial^2 \xi / \partial t^2 = v^2 \, \partial^2 \xi / \partial x^2$, находим \textbf{скорость продольной волны в тонком стержне}:
\begin{equation}
    v_{\parallel} = \sqrt{\frac{E}{\rho}}. \label{eq:rod_speed}
\end{equation}
\end{theorem}

\begin{definition}
\textbf{Модуль Юнга} $E$ --- физическая величина, характеризующая способность материала сопротивляться растяжению или сжатию при упругой деформации. Определяется как отношение механического напряжения к относительной деформации:
\begin{equation}
    E = \frac{\sigma}{\varepsilon} = \frac{F/S}{\Delta l / l}.
\end{equation}
Единица измерения --- паскаль (Па). Чем больше $E$, тем жёстче материал. Например: сталь $E \approx 2 \cdot 10^{11}$ Па, алюминий $E \approx 7 \cdot 10^{10}$ Па.
\end{definition}

\begin{property}[Скорость поперечных волн в изотропной твёрдой среде]
В изотропном твёрдом теле могут распространяться как продольные, так и поперечные волны. Для \textbf{поперечных волн} (деформация сдвига) скорость определяется модулем сдвига $G$ (модуль жёсткости):
\begin{equation}
    v_{\perp} = \sqrt{\frac{G}{\rho}}. \label{eq:shear_speed}
\end{equation}
Для большинства материалов $G < E$, поэтому поперечные волны распространяются медленнее продольных: $v_{\perp} < v_{\parallel}$. В жидкостях и газах $G = 0$, поэтому поперечные упругие волны в них не распространяются.
\end{property}

\begin{remark}
\textbf{Сравнение скоростей:}
\begin{itemize}
    \item Продольные волны в стержне: $v_{\parallel} = \sqrt{E/\rho}$;
    \item Поперечные волны в твёрдом теле: $v_{\perp} = \sqrt{G/\rho}$;
    \item Звук в газе: $v = \sqrt{\gamma P / \rho} = \sqrt{\gamma RT / M}$.
\end{itemize}
Для стали ($E \approx 2 \cdot 10^{11}$ Па, $\rho \approx 7800$ кг/м³): $v_{\parallel} \approx 5000$ м/с. Это значительно выше скорости звука в воздухе ($\approx 340$ м/с), что объясняет, почему звук по рельсу слышен раньше, чем по воздуху.
\end{remark}

\begin{figure}[htbp]
    \centering

    % === РИСУНОК 1: Элемент стержня ===
    \begin{tikzpicture}[scale=1.2, >=Stealth, font=\small]
        \node[font=\bfseries, align=center] at (1.5, 1.8) {Элемент стержня $dx$};

        % Стержень
        \fill[blue!10] (0, 0) rectangle (3, 1);
        \draw[thick] (0, 0) rectangle (3, 1);

        % Силы (Красные)
        \draw[red, ->, thick] (-1, 0.5) -- (0, 0.5) node[midway, above] {$\sigma(x)$};
        \draw[red, ->, thick] (4.5, 0.5) -- (3, 0.5) node[midway, above] {$\sigma(x+dx)$};

        % Смещения (Синие)
        \draw[blue, ->, thick] (0.1, 0.5) -- (0.6, 0.5) node[midway, above] {$\xi(x)$};
        \draw[blue, ->, thick] (1.5, 0.5) -- (3, 0.5) node[midway, above] {$\xi(x+dx)$};

        % Размеры
        \draw[<->] (0, -0.5) -- (3, -0.5) node[midway, fill=white] {$dx$};
        \node[right] at (3, 1) {$S$};

        % Формула деформации
        \node[draw, fill=white, align=center] at (1.5, -1.2) {$\varepsilon = \frac{\partial \xi}{\partial x}$};
    \end{tikzpicture}

    \vspace{1.5cm} % Большой отступ между рисунками

    % === РИСУНОК 2: Типы волн ===
    \begin{tikzpicture}[scale=1.2, >=Stealth, font=\small]
        \node[font=\bfseries] at (0, 2.2) {Продольная ($v_1$)};

        % Блок
        \fill[gray!20] (-1, 0.5) rectangle (1, 1.5);
        \draw[thick] (-1, 0.5) rectangle (1, 1.5);

        % Колебания (Вдоль)
        \foreach \y in {0.7, 1.0, 1.3} {
            \draw[red, <->] (-0.8, \y) -- (-0.2, \y);
        }

        % Направление
        \draw[->, thick] (1.2, 1) -- (2.5, 1) node[right] {$\vec{v}_1$};
        \node[align=center, font=\tiny] at (0, -0.2) {$v_1 = \sqrt{\frac{E}{\rho}}$};

        % Разделитель
        \draw[gray, thick] (3.5, 0) -- (3.5, 2);

        % Поперечная
        \node[font=\bfseries] at (5.5, 2.2) {Поперечная ($v_\perp$)};

        % Блок
        \fill[gray!20] (4.5, 0.5) rectangle (6.5, 1.5);
        \draw[thick] (4.5, 0.5) rectangle (6.5, 1.5);

        % Колебания (Поперек)
        \foreach \x in {4.7, 5.1, 5.5, 5.9, 6.3} {
            \draw[red, <->] (\x, 0.7) -- (\x, 1.3);
        }

        % Направление
        \draw[->, thick] (6.7, 1) -- (8, 1) node[right] {$\vec{v}_\perp$};
        \node[align=center, font=\tiny] at (5.5, -0.2) {$v_\perp = \sqrt{\frac{G}{\rho}}$};

        % Сноска
        \node[draw, fill=yellow!10, font=\tiny] at (3.5, -0.8) {В жидкостях $G=0 \Rightarrow v_\perp=0$};
    \end{tikzpicture}

    \caption{Слева: силы упругости на гранях элемента стержня. Справа: направления колебаний частиц для продольных и поперечных волн.}
    \label{fig:stresses_and_waves}
\end{figure}

